%%%%%%%%%%%%%%%%%%%%%%%%%%%%%%%%%%%%%%%%%%%%%%%%%%%%%%%%%%%%
\section{形式化分析}
\label{sec:formal-analysis}

我们证明认证工作流在假设 L1-L3 下针对攻击者 A1-A5 实现了 O1-O5，通过围绕意图和完整性组织的七个引理。完整性属性（引理1、2、3）确保操作是真实的且防篡改的。意图属性（引理4、5）确保操作满足策略并安全组合。完备性属性（引理6、7）确保所有操作都被验证，且验证是独立的。

\noindent\textbf{主要安全定理。} \textbf{定理（认证工作流安全性）}：在信任假设 L1-L3 下，认证工作流针对具有 A1-A5 能力的攻击者实现了完整性（O1）、策略执行（O2）、权限非升级（O3）、上下文完整性（O4）和可问责性（O5）。

\noindent\textbf{完整性属性。} 这些引理确保操作是真实的、防篡改的且不可否认的——针对具有组件妥协（A3）和网络攻击（A4）能力的攻击者。

\textbf{引理1（真实性）}：所有被接受的调用都具有有效的密码学签名，将主体绑定到操作、参数、策略和上下文状态。

\textit{证明}：策略执行点（Policy Enforcement Point，PEP）使用从智能体注册表（第~\ref{sec:enforcement}~节）获取的主体公钥验证签名。在密码学困难性假设 L1 下，没有私钥而伪造签名在计算上是不可行的——即使具有应用控制（A1）、网络访问（A4）或组件妥协（A3）能力的攻击者也无法为其他主体的身份生成有效签名。这实现了完整性（O1）：操作在密码学上是真实的。$\square$

\textbf{引理2（防篡改）}：对上下文状态或审计日志的修改在密码学上是可检测的，且重放攻击被阻止。

\textit{证明}：上下文状态通过链接顺序状态的哈希链维护：$h_i = \text{Hash}(h_{i-1} \parallel \text{state}_i \parallel \text{sequence}_i)$。每个调用包含一个绑定在签名中的单调递增序列号。修改任何状态都需要找到哈希碰撞；重放旧调用会被检测到，因为 PEP 拒绝序列号小于或等于先前已处理值的调用。在密码学困难性（L1）下，找到碰撞在计算上是不可行的。多轮操纵攻击（A5）试图渐进式状态投毒或重放操作会被检测到，因为篡改会破坏哈希链或违反序列单调性。审计日志使用相同的哈希链机制。这实现了上下文完整性（O4）。$\square$

\textbf{引理3（不可否认性）}：主体不能否认执行了审计日志中记录带有有效签名的操作。

\textit{证明}：只有持有私钥的主体才能生成有效签名（引理1）。审计日志维护防篡改记录（引理2）。因此，审计日志中的已签名操作构成了主体行为的密码学证明。这实现了可问责性（O5）。$\square$

\noindent\textbf{意图属性。} 这些引理确保操作满足策略并安全组合——针对具有应用控制（A1）、内容注入（A2）和渐进式攻击（A5）能力的攻击者。

\textbf{引理4（策略执行）}：所有已执行操作满足由组织层级和资源约束组合而成的有效策略，且策略标识符在密码学上绑定在签名中。

\textit{证明}：每个调用在签名载荷中包含一个策略标识符。PEP 检索策略并通过所有适用策略的交集计算有效策略 $P_{\text{eff}}$——组织策略（基础、部门、团队）和资源特定策略（第~\ref{sec:policy}~节）。试图进行策略替换攻击的攻击者（能力 A1、A4）会失败，因为修改策略标识符会使签名失效（引理1）。策略交集确保单调约束（第~\ref{sec:policy}~节，定理1-3）——组合策略只会变得更严格，绝不会更宽松。这实现了策略执行（O2）和权限非升级（O3）。$\square$

\textbf{引理5（组合安全）}：当工作流跨越多个框架时，安全属性在框架边界之间得到保留。

\textit{证明}：考虑框架 F1 上的操作 O1 调用框架 F2 上的操作 O2。两个调用都携带绑定主体和策略的签名（引理1）。框架 F2 的 PEP 独立验证 O2 的签名，并将有效策略计算为 F1 策略和 F2 策略的交集（引理4）。框架 F1 无法绕过 F2 的验证，因为：(a) F2 的 PEP 独立验证签名（引理7，见下文证明），且 (b) 策略交集确保 O2 必须同时满足 F1 和 F2 的约束。跨 N 个框架的组合创建了 N 个独立验证点，组合策略为 $P_{\text{eff}} = P_1 \cap P_2 \cap \ldots \cap P_N$，实现了单调约束。第~\ref{sec:enforcement}~节在九个框架组合上演示了这一点。$\square$

\noindent\textbf{完备性属性。} 这些引理确保所有操作都被验证，且验证在各个控制面之间是独立的。

\textbf{引理6（控制面完备性与极小性）}：四个控制面 \{S1, S2, S3, S4\} 是完备的（所有资源访问操作至少跨越一个面）且极小的（每个面都是必要的）。

\textit{通过枚举与必要性证明}：

\textit{完备性：}我们枚举所有资源访问操作，并展示每个操作至少跨越一个面：(1) \textit{计算资源}——LLM 推理跨越 S1（提示携带指令）；工具执行跨越 S2（特权操作）。(2) \textit{数据资源}——RAG 检索、数据库查询、网页抓取跨越 S3（外部数据流入推理）——被内容注入攻击（A2）利用。(3) \textit{状态资源}——会话状态、工作流上下文、记忆跨越 S4（跨轮次的持久状态）——被多轮操纵（A5）针对。(4) \textit{跨智能体通信}——委托跨越 S1（指令传播）或 S2（调用）加上 S4（上下文继承）。不访问外部资源的操作（纯计算）不需要保护——它们无法泄露数据或违反策略。

\textit{极小性：}每个面都是必要的——移除任何面都会使攻击失去防护：移除 S1（提示）：通过 S3 数据$\rightarrow$LLM 的间接提示注入绕过验证。移除 S2（工具）：尽管有提示验证，仍可进行未授权工具执行。移除 S3（数据）：尽管有 S1/S2 验证，被投毒的 RAG 数据仍可导致工具调用。移除 S4（上下文）：跨多轮工作流的会话劫持和上下文投毒。因此，\{S1, S2, S3, S4\} 是完备且极小的。第~\ref{sec:enforcement}~节在跨越三个架构层的九个框架上对此进行了实证验证。$\square$

\textbf{引理7（PEP 独立性）}：攻陷一个策略执行点（PEP）不会削弱其他 PEP 的验证。

\textit{证明}：每个 PEP 仅使用以下要素独立验证操作：(1) 用于签名和哈希链验证的密码学原语（L1），(2) 从可信控制面检索的公钥和策略（L2），以及 (3) 其自身的验证逻辑（L3）。PEP 之间不共享运行时状态，也不进行协调。攻陷某个控制面上的 PEP（能力 A3）的攻击者只能绕过该控制面的验证——其他控制面使用各自的 PEP 实例执行独立验证。多面攻击需要攻陷多个独立的 PEP，每个都受 L3 保护。这实现了纵深防御（defense-in-depth）：爆炸半径仅限于被攻陷 PEP 所在的控制面。$\square$

\noindent\textbf{主定理证明。} \textit{主定理证明}：根据引理6，所有资源访问至少跨越一个控制面。每个控制面内嵌一个 PEP 执行独立验证（引理1-4、7）。多框架工作流在框架边界之间组合验证（引理5）。验证通过策略评估执行意图（引理4），通过签名（引理1）和哈希链（引理2）执行完整性。密码学证言（attestation）通过已签名、哈希链式声明提供工作流依赖关系。

每个资源访问操作有两种结果之一：(1) \textit{授权执行}——操作满足所有密码学检查（有效签名、防篡改上下文、有效序列号）和策略约束（允许的资源、参数在范围内、所需的密码学证言存在），执行后以不可否认的审计追踪记录（引理3）；(2) \textit{被阻止}——操作至少未通过一项检查（无效签名、策略违反、上下文篡改、序列号违规或缺失密码学证言），在执行前被拒绝并记录违规。

\textit{具体示例}：在 Q4 攻击（第~\ref{sec:problem}~节）中，被投毒的数据指示 LLM 泄露凭证。当 LLM 生成对 \texttt{credentials.db} 的文件系统读取操作时，工具的 PEP 评估策略——该策略拒绝匹配 \texttt{*credential*} 的路径——从而阻止该操作，尽管来自已认证 LLM 的签名是有效的。策略执行防止未授权访问，无论签名是否有效。即使多个边界被攻陷，每个控制面执行独立验证（引理7），限制了爆炸半径。

因此，在 L1-L3 下，具有 A1-A5 能力的攻击者无法在不受检测的情况下执行未授权资源访问，从而实现了 O1-O5。$\square$

\noindent\textbf{实践启示。} 这些形式化属性支撑了第~\ref{sec:trust-layer}~节中的操作性解决方案：引理1使能主体传播，引理4+5使能可验证委托，引理2使能会话级前向安全，引理3使能双向可问责性。

